\section{光的干涉}
\sourceproblem{选择1}{41}{47}
\problemtext 真空波长 $\lambda$ 的光在介质中传播，相位差为 $\pi$，求光程。
\workings $\Delta\phi=2\pi ns/\lambda$，故光程 $ns=\lambda/2$，选A。
\crosscheck 若只给模 $2\pi$ 的相位差，则光程还可为 $(j+1/2)\lambda$；教材取实际累积相位差。
\sourceproblem{选择2}{41}{47}
\problemtext 正入射薄膜，上、膜、下介质折射率依次为 $n_1<n_2<n_3$，膜厚 $e$，比较反射光光程。
\workings 两次反射均由低折射率射向高折射率，均有 $\pi$ 相变，反射相变差为零。额外传播光程为 $2n_2e$，选A。
\crosscheck 干涉相位应以真空波长计算；两束均有半波损失时不能再单独添加半波。
\sourceproblem{选择3}{41}{47}
\problemtext 双缝各覆盖相同厚度 $e$、不同折射率 $n_1,n_2$ 的介质片，求原中央极大处光程差。
\workings 原两路径几何等长，介质引入的差为 $(n_2-n_1)e$（第二路减第一路），选B。
\crosscheck 反向定义则变号；相同折射率时为零。
\sourceproblem{选择4}{41}{47}
\problemtext 双缝板略向下平移，其他条件保持，判断屏上条纹。
\workings 在正入射平行照明和近轴模型下，中央明纹随双缝中点下移，条纹间距 $\lambda D/a$ 不变，选A。
\crosscheck 若实际是有限距离固定点源，入射相位也会改变；采用教材标准双缝模型。
\sourceproblem{选择5}{41}{47}
\problemtext 空气劈尖左薄右厚，图示反射干涉条纹向右弯曲一个条纹间距，判断工件缺陷。
\workings 同一条纹对应同一空气层厚度。右移本应使空气层增厚，故工件局部凸起补偿该厚度，凸起高度为相邻同类条纹厚度差 $\lambda/2$，选B。
\crosscheck 依据图2弯曲顶点到下一条纹直线的距离为一整条纹间距，不是半间距。
\sourceproblem{选择6}{41}{47}
\problemtext 牛顿环装置平凸透镜缓慢下降接近平板，判断环纹运动。
\workings 空气隙 $h=h_0+r^2/(2R)$。同级条纹的 $h$ 固定，$h_0$ 减小使 $r$ 增大，故向外扩张，选C。
\crosscheck 这里为接触前平移的理想球面模型；接触后弹性形变需另分析。
\sourceproblem{选择7}{42}{48}
\problemtext 迈克耳孙干涉仪反射镜移动距离 $d$，求经过参考点的条纹数。
\workings 光往返镜面，光程变化 $2d$，条纹移动量 $2d/\lambda$，选B。
\crosscheck 若不是整周期，表达式为移动的条纹周期数；实际完整通过的整数条数还与初相有关。
\sourceproblem{填空1}{42}{48}
\problemtext 同样时间内单色光在空气和玻璃中传播的路程是否相等？
\workings 在通常玻璃折射率大于空气的模型中，传播速率 $c/n$ 不同，因此几何路程不相等；光程 $ns=ct$ 相等（相速度模型）。
\crosscheck 区分几何距离与光程；在色散介质中脉冲群速度须另行处理。
\sourceproblem{填空2}{42}{48}
\problemtext 两相干源至P距离 $r_1,r_2$，分别垂直穿过厚度 $t_1,t_2$、折射率 $n_1,n_2$ 的介质板，求光程差。
\workings 第二路减第一路为 $\Delta=r_2-r_1+(n_2-1)t_2-(n_1-1)t_1$。
\crosscheck 其余路段为真空；若两板均去掉，恢复几何路程差。
\sourceproblem{填空3}{42}{48}
\problemtext 等距同相源的一路插入厚度 $e$、折射率 $n$ 的玻璃片，求相位差；波长500 nm、$n=1.5$，P为第三级明纹时求厚度。
\workings 插片一路多出光程 $(n-1)e$，相位差 $2\pi(n-1)e/\lambda$。三级明纹要求 $(n-1)e=3\lambda$，故 $e=3000$ nm。
\crosscheck 以零级中央明纹为级次起点；反向定义相位差则变号。
\sourceproblem{填空4}{42}{48}
\problemtext 双缝明纹处遮住第二缝，沿两缝中垂面放高折射率反射面，判断原明纹点。
\workings 镜面形成第一缝的虚像，几何位置对应第二缝；反射多出 $\pi$ 相变，原整数波长光程差转为半整数波长的相消条件，故为暗纹。
\crosscheck 理想完全暗纹还需两束振幅相等；不等时是干涉极小。
\sourceproblem{填空5}{42}{48}
\problemtext 空气中厚 $e$、折射率 $n>1$ 的薄膜正入射反射干涉光程差。
\workings 上表面反射有半波相变、下表面无，等效光程差为 $2ne+\lambda/2$（模 $\lambda$）。
\crosscheck 写成 $2ne-\lambda/2$ 亦等价；两式相差整波长。
\sourceproblem{填空6}{42}{48}
\problemtext 双缝在空气中条纹间距1.5 mm，整体浸入折射率1.33的水中求间距。
\workings $\Delta x=\lambda_0D/(na)$，因此 $\Delta x'=1.5/1.33\approx1.13$ mm。
\crosscheck 默认几何尺寸不变且整体浸没，频率不变。
\sourceproblem{填空7}{42}{48}
\problemtext 空气中折射率1.33的薄膜，580 nm光垂直入射，求反射光极小的最小非零厚度。
\workings 一次反射有半波相变，反射相消要求 $2ne=j\lambda$。取最小正整数 $j=1$，$e_{\min}=580/(2\times1.33)\approx218$ nm。
\crosscheck 零厚度为退化情形；完全相消还要求反射振幅匹配。
\sourceproblem{填空8}{42}{48}
\problemtext 折射率 $n$ 的劈尖膜正入射，求第三与第六条明纹的厚度差。
\workings 相邻明纹的厚度差 $\lambda/(2n)$，跨三间隔得 $3\lambda/(2n)$。
\crosscheck 不依赖明纹起始级次的约定。
\sourceproblem{填空9}{42}{48}
\problemtext 空气劈尖左为棱边，上玻璃板竖直上移而倾角不变，判断条纹间隔及移动方向。
\workings $h=h_0+x\tan\theta$，同级厚度固定。$h_0$ 增大使 $x$ 减小，条纹向近棱边方向移；间距 $\lambda/(2\tan\theta)$ 不变。
\crosscheck 与改变劈尖夹角的情形不同。
\sourceproblem{填空10}{43}{49}
\problemtext 折射率1.52玻璃组成牛顿环，由空气移入折射率1.33的水中，判断环纹疏密。
\workings 反射暗环 $r_j^2=j\lambda_0R/n$，各级半径及相邻径向间距缩为原来的 $1/\sqrt{1.33}$，故变密。
\crosscheck 水折射率仍低于玻璃，反射相变关系不变。
\sourceproblem{填空11}{43}{49}
\problemtext 迈克耳孙干涉仪一路插入折射率 $n$ 的透明膜，光程差改变一波长，求厚度。
\workings 光两次穿膜，$2(n-1)e=\lambda$，故 $e=\lambda/[2(n-1)]$。
\crosscheck 假定法向穿膜且周围为空气或真空近似；单程干涉仪不含因子2。
\sourceproblem{计算1}{43}{49}
\problemtext 双缝波长600 nm、缝距 $3\times10^{-4}$ m、屏距1 m，求两侧五级明纹间距；一缝加厚 $5.5\times10^{-6}$ m、折射率1.33玻璃片，求零级移至何处。
\workings 近轴间距 $\Delta x=\lambda D/a=2.00$ mm，两侧五级中心相距 $10\Delta x=20.0$ mm。附加光程 $(n-1)e=1.815\times10^{-6}$ m，零级朝覆盖玻璃片的缝侧移动，
\[
\Delta j=\frac{(n-1)e}{\lambda}=3.025,\qquad
|\Delta x_0|=6.05\,\mathrm{mm}.
\]
即靠近原第三级明纹，但严格并非恰在原整数级明纹中心。
\crosscheck 原PDF49的厚度及折射率已读清；若按有效数字近似可称第三级，不能把3.025硬写成精确3。近轴条件 $|x|\ll D$ 满足。

\sourceproblem{计算2}{44}{50}
\problemtext 白光正入射空气中的肥皂膜，厚380 nm，折射率1.33，分析反射颜色。
\workings 两次反射只有一次发生半波相变，反射增强条件 $2ne=(j+1/2)\lambda$，其中 $2ne=1010.8$ nm。可见光范围内主要增强波长为 $1010.8/1.5=673.9$ nm；下一增强波长404.3 nm位于可见紫端。反射减弱条件 $2ne=j\lambda$，可见范围内约505.4 nm的绿光减弱。因此理想白光照明下呈偏红紫色，实际色调取决于光源谱及视觉响应。
\crosscheck 不把所有可见光简化为唯一波长；薄膜反射率有限，增强和减弱是光谱调制。
\sourceproblem{计算3}{44}{50}
\problemtext 硅表面的二氧化硅薄膜磨成劈形，折射率分别为3.42、1.50；以550 nm光正入射，AB段见6条暗纹且B恰为暗纹，求原膜厚。
\workings 空气至膜、膜至硅的两次反射均有半波相变，相对相变为零。因此反射暗纹满足 $2ne=(j+1/2)\lambda$。A处膜厚趋零，为亮纹而非暗纹，六条暗纹对应 $j=0,1,\ldots,5$，B处膜厚
\[
e=\frac{(5+1/2)550}{2(1.50)}\,\mathrm{nm}\approx1008.3\,\mathrm{nm}.
\]
\crosscheck 原图A处为磨至基底的尖端；若实际起始厚度非零，须另确定级次。相邻暗纹厚差为183.3 nm。
\sourceproblem{计算4}{45}{51}
\problemtext 同一牛顿环装置分别用已知波长 $\lambda_1$ 与未知波长 $\lambda_2$ 照明，测第四和第九暗环的半径差为 $l_1,l_2$，求未知波长。
\workings 反射暗环满足 $r_j=\sqrt{jR\lambda}$（接触、空气隙）。因此 $r_9-r_4=(3-2)\sqrt{R\lambda}$，两次测量相除得
\[
\lambda_2=\lambda_1(l_2/l_1)^2.
\]
\crosscheck 必须平方半径差的比；若存在初始非零空气隙，不能直接使用同一简化式。
